\section{Backend: Configuración Inicial y Boilerplate}
\label{sec:s0_backend}

El andamiaje del backend (\comando{ethoshubB}) se construyó sobre Spring Boot 4.0.6 con Maven,
priorizando la reproducibilidad del entorno y la integración temprana del flujo de
empaquetado monolítico. Esta sección documenta la inicialización del repositorio
(Sección~\ref{ssec:s0_be_init}) y la estructura de directorios junto con el flujo de CI/CD
(Sección~\ref{ssec:s0_be_estructura}).\par

\subsection{Inicialización del Entorno y Gestión de Repositorios}
\label{ssec:s0_be_init}

Se adoptó una estrategia de ramas \textbf{GitFlow} (\comando{main}, \comando{develop},
\comando{feature/*}, \comando{release/*}), con \textit{tags} semánticos por sprint
(\comando{release/v0.1.0} en adelante). El proyecto se generó con la estructura estándar de
Spring Boot y el \textit{wrapper} de Maven (\comando{mvnw}) para garantizar la paridad de la
versión de Maven entre desarrolladores. El \comando{pom.xml} declara las dependencias núcleo
fijadas en este sprint:

\begin{itemize}[noitemsep]
    \item \comando{spring-boot-starter-web} --- API REST con Spring Web MVC.
    \item \comando{spring-boot-starter-security} + \comando{oauth2-resource-server} --- validación de JWT.
    \item \comando{spring-boot-starter-validation} --- Jakarta Bean Validation en los DTO.
    \item \comando{spring-boot-starter-jooq} / jOOQ --- acceso a datos tipado contra \comando{core}.
    \item \comando{spring-boot-starter-actuator} --- \textit{health checks} y métricas.
    \item \comando{springdoc-openapi} --- generación de Swagger UI.
\end{itemize}

El archivo \comando{application.yml} se parametriza con \textit{placeholders}
\comando{\$\{VAR:default\}} para externalizar la configuración (URL de Supabase, claves, SMTP),
sentando la base de la gestión de secretos que se formaliza en el Sprint 5
(Sección~\ref{sec:s5_backend}).

\begin{lstlisting}[language=Java, caption={Punto de entrada del backend (Sprint 0).}, label={lst:s0_main}]
@SpringBootApplication
public class EthosHubApplication {
    public static void main(String[] args) {
        SpringApplication.run(EthosHubApplication.class, args);
    }
}
\end{lstlisting}

\subsection{Estructura de Directorios, Enrutamiento Base y CI/CD}
\label{ssec:s0_be_estructura}

La estructura del proyecto refleja el patrón \textit{vertical slice} (un paquete por dominio)
fijado en la Sección~\ref{sec:s0_arquitectura}. El árbol base entregado en el Sprint 0 es el
siguiente:

\begin{lstlisting}[caption={Estructura de directorios del backend.}, label={lst:s0_tree}]
ethoshubB/
|-- build.sh                  # Empaquetado monolitico (FE + BE -> JAR)
|-- pom.xml                   # Dependencias y plugins Maven
|-- Dockerfile                # Imagen de despliegue
`-- src/
    |-- main/
    |   |-- java/com/ethoshub/
    |   |   |-- EthosHubApplication.java
    |   |   |-- auth/ security/ profile/ project/ ...
    |   |   |-- config/        # Seguridad, CORS, filtros
    |   |   `-- shared/        # ApiResponse, excepciones, util
    |   `-- resources/
    |       |-- application.yml
    |       |-- db/migration/V*.sql
    |       `-- static/        # dist del frontend (generado)
    `-- test/java/             # JUnit 5 + Testcontainers
\end{lstlisting}

El \textbf{enrutamiento base} se organiza por prefijos versionados (\comando{/api/v1/...}),
con un \comando{HealthController} y la respuesta envuelta de forma homogénea por
\comando{shared/ApiResponse}. La configuración de seguridad inicial (\comando{config/}) deja
preparados los filtros que el Sprint 1 implementa por completo
(Sección~\ref{sec:s1_backend}).

Respecto al \textbf{flujo de CI/CD}, el directorio \comando{.github/} aloja los \textit{workflows}
de integración continua, y el \textit{script} \comando{build.sh} encapsula el empaquetado
monolítico. La tabla~\ref{tab:s0_buildsh} documenta sus modos de operación.

\begin{table}[!ht]
    \centering
    \renewcommand{\arraystretch}{1.35}
    \fontsize{9}{10.8}\selectfont\sffamily
    \begin{tabularx}{\textwidth}{|l|X|}
        \hline
        \rowcolor{colorPrimario}
        \textcolor{white}{\textbf{Comando}} & \textcolor{white}{\textbf{Acción}} \\ \hline
        \comando{./build.sh} & \textit{Build} completo: compila el FE (\comando{npm ci \&\& build}) y genera el JAR. \\ \hline
        \rowcolor{grisTecnico}
        \comando{./build.sh --skip-fe} & JAR solo backend, reutilizando el \comando{dist/} existente. \\ \hline
        \comando{./build.sh --prod} & Activa el perfil Spring \comando{prod} (\comando{application-prod.yml}). \\ \hline
    \end{tabularx}
    \caption{Modos de operación del \textit{script} \comando{build.sh}.}
    \label{tab:s0_buildsh}
\end{table}

\begin{NotaTecnica}
El flujo de \comando{build.sh} compila \comando{../ethoshubF} con Node, copia \comando{dist/} a
\comando{src/main/resources/static/} mediante el \comando{maven-resources-plugin} (fase
\comando{process-resources}) y empaqueta el JAR con el \comando{spring-boot-maven-plugin}. Este
mecanismo es la columna vertebral del despliegue monolítico documentado en el Sprint 5.
\end{NotaTecnica}

\subsubsection{Etapas del Flujo de Integración Continua}
La tabla~\ref{tab:s0_cicd} resume las etapas del \textit{pipeline} de CI definido en
\comando{.github/}, ejecutadas en cada \textit{push} a las ramas de integración.

\begin{table}[!ht]
    \centering
    \renewcommand{\arraystretch}{1.3}
    \fontsize{9}{10.8}\selectfont\sffamily
    \begin{tabularx}{\textwidth}{|l|X|}
        \hline
        \rowcolor{colorPrimario}
        \textcolor{white}{\textbf{Etapa}} & \textcolor{white}{\textbf{Acción}} \\ \hline
        Checkout & Clona el repositorio y configura JDK 17. \\ \hline
        \rowcolor{grisTecnico}
        Build & \comando{mvn clean package} (compila y empaqueta). \\ \hline
        Test & JUnit 5 + Testcontainers (PostgreSQL). \\ \hline
        \rowcolor{grisTecnico}
        Package FE & \comando{npm ci \&\& build} y copia a \comando{static/}. \\ \hline
        Artefacto & Publica el JAR monolítico versionado. \\ \hline
    \end{tabularx}
    \caption{Etapas del flujo de integración continua (Sprint 0).}
    \label{tab:s0_cicd}
\end{table}

\subsection{Diagrama de Secuencia: Empaquetado Monolítico (\texttt{build.sh})}
\label{ssec:s0_seq_build}
La Figura~\ref{fig:s0_seq} muestra la secuencia del empaquetado FE+BE en un único JAR.

\begin{figure}[!ht]
    \centering
    \begin{tikzpicture}[font=\sffamily\scriptsize, >=Latex,
        part/.style={draw=colorPrimario, fill=headerTabla, rounded corners=2pt,
            minimum height=0.6cm, text width=2.1cm, align=center},
        msg/.style={-{Latex[length=1.6mm]}, colorPrimario},
        ret/.style={-{Latex[length=1.6mm]}, dashed, grisISO}]
        \node[part] (d)  at (0,0)    {Desarrollador};
        \node[part] (b)  at (3.6,0)  {build.sh};
        \node[part] (n)  at (7.4,0)  {Node / Vite};
        \node[part] (m)  at (11.2,0) {Maven};
        \foreach \n in {d,b,n,m} \draw[grisISO, dashed] (\n.south) -- ++(0,-4.6);
        \draw[msg] (0,-0.9)    -- node[above]{1. ./build.sh --prod} (3.6,-0.9);
        \draw[msg] (3.6,-1.7)  -- node[above]{2. npm ci \&\& build} (7.4,-1.7);
        \draw[ret] (7.4,-2.5)  -- node[above]{3. dist/} (3.6,-2.5);
        \draw[msg] (3.6,-3.3)  -- node[above]{4. package (dist$\to$static)} (11.2,-3.3);
        \draw[ret] (11.2,-4.1) -- node[above]{5. JAR monolítico} (3.6,-4.1);
    \end{tikzpicture}
    \caption{Diagrama de secuencia del empaquetado monolítico con \comando{build.sh}.}
    \label{fig:s0_seq}
\end{figure}

\clearpage
